\chapter{讨论}\label{chap:discussion}

\section{主要结果的含义}

\subsection{RBSOG当前尚未体现性能优势}

本文最清晰的结果是：在当前纯Python/JIT混合原型和本文可运行规模内，RBSOG显著慢于PPPM。主基准中，RBSOG平均步时间约为PPPM的\RBSOGSlowdownVsPPPM 倍；规模扫描中，慢速比从 $N=320$ 时的21.982扩大到 $N=1280$ 时的\ScalingSlowdownAtMaxN。稳健统计也支持这一点：步时间bootstrap比例区间为 $[\BootstrapStepRatioLow,\BootstrapStepRatioHigh]$，没有接近1。

这一结果并不否定RBSOG路线的长期动机。RBSOG希望减少对全局FFT通信路径的依赖，而当前实验是在单机串行原型上完成的，并未进入其潜在优势最可能体现的大规模并行环境。然而，从毕业论文实证角度看，当前数据必须被如实表述为：RBSOG在现有实现下还没有转化为实测性能优势。

\subsection{压力方差结论比步时间结论更不确定}

压力方差均值上，RBSOG高于PPPM，主基准比例为\BootstrapVarianceRatio。但bootstrap区间为 $[\BootstrapVarianceRatioLow,\BootstrapVarianceRatioHigh]$，跨过1，置换检验也不显著。这说明压力方差结论不能像步时间结论那样强。更准确的表述是：当前默认参数下，RBSOG没有表现出稳定的压力方差优势，并且若干种子存在明显长尾波动。
这种差异符合随机批量方法的基本特征。步时间主要由实现路径和样本数量决定，跨种子变化较小；压力方差则同时受采样噪声、初始条件、温控压控反馈和有限轨迹长度影响，因此更容易出现高不确定性。后续如果要对压力稳定性做更强结论，需要更长轨迹、更大种子数和更丰富的统计检验。

\subsection{结构代理量与压力方差并不等价}

长轨迹实验给出一个值得注意的现象：RBSOG压力方差显著高于PPPM，但面积/脂质和膜厚代理量的标准差、漂移率在5000步统计下并不更差。这说明结构代理量和压力方差不能互相替代。压力是对瞬时力、virial和采样噪声更敏感的全局量；面积/脂质和膜厚代理量则在Berendsen控制和简化代理体系下更平滑。
因此，本文不把“压力方差更高”直接等同于“结构一定更差”，也不把“结构代理量更平滑”直接等同于“算法更稳定”。更稳妥的结论是：当前RBSOG的统计压力行为仍需控制，而结构代理量没有给出单向否定或单向肯定的证据。

\section{误差来源与参数敏感性}

\subsection{批量大小是有效但昂贵的误差控制旋钮}

direct reference实验显示，增大批量大小 $P$ 通常能降低相对力RMSE。例如在 $N=640$ 体系上，$P=25$ 的力RMSE为1.762，$P=400$ 降至0.942。这说明随机批量采样的误差确实可以通过增加样本量缓解。

问题在于，这种缓解非常昂贵。批量扫描中，$P=400$ 的压力方差最低，但步时间最高；$P=50$ 最快，但方差更高。$P=\RecommendedBatch$ 是当前加权目标下的折中点，而不是全局支配点。这种结果说明RBSOG目前主要处在“可调但代价高”的阶段。

\subsection{SOG项数不是单调收益参数}

另一个重要发现是，增加SOG项数 $M$ 并不一定降低误差。在当前实现中，$M=6$ 往往优于 $M=12,18,24$。这说明当前高斯和核参数生成、最小二乘拟合区间、非负权重裁剪等细节对误差有很大影响。
从方法学角度看，这一点很关键。如果把RBSOG的误差简单归因于“随机批量采样噪声”，就会忽略SOG核拟合本身的误差贡献。后续工作应把误差分解为至少两部分：一是固定核下的随机采样误差，二是不同SOG核参数化方式造成的系统误差。

\subsection{PPPM公平性基线的边界}

本文通过网格公平性扫描选择 $\ChosenPPPMGrid^3$ 作为PPPM基线。这有助于避免把过高成本的PPPM配置直接拿来与RBSOG比较，但也带来一个限制：该PPPM配置不是高精度PME/PPPM生产参数。direct reference结果显示，公平性选择下的PPPM力RMSE也较高。
因此，本文所有“RBSOG与PPPM”的结论都应理解为当前原型公平配置下的比较，而不是对成熟软件中高精度PPPM的最终评价。这个限定不会削弱“RBSOG当前更慢”的结论，反而使其更严格：即便在较低成本PPPM基线下，RBSOG仍未取得性能优势。

\section{Profiling揭示的优化方向}

\subsection{当前瓶颈集中在长程采样路径}

Profiling结果显示，PPPM耗时主要集中在FFT求解，占\PPPMFFTShare\%；RBSOG耗时主要集中在随机采样、长程核计算和样本过滤，三者合计约占 $\RBSOGSamplingShare + \RBSOGLongRangeShare + \RBSOGFilterShare$ 的量级。短程近邻表不是主要瓶颈。
这意味着后续优化不应优先围绕短程pair list做小修小补，而应集中处理长程采样路径：减少随机索引生成开销，避免生成后再大量过滤无效样本，改进内存布局，并将长程核改写为更适合向量化、并行化或GPU执行的数据流。

\subsection{JIT说明实现层仍有明显空间}

JIT消融显示，开启JIT带来\JITSpeedup 倍加速，而压力方差不变。这说明当前RBSOG慢速中确实存在大量实现层开销。与此同时，即便开启JIT，RBSOG仍明显慢于PPPM，因此JIT只能证明“还有工程优化空间”，不能证明“工程优化后必然超过PPPM”。
更现实的路线是分阶段推进：先把采样和过滤路径降到合理成本，再比较CPU多线程或GPU版本中的相对扩展性，最后才讨论是否能够在大规模并行场景中规避FFT通信瓶颈。

\section{与相关工作的关系}

传统Ewald、PME和PPPM方法通过实空间与倒易空间分解，把长程静电计算复杂度降到适合大规模MD的水平\cite{darden1993particle,hockney1988computer}。它们的短板主要在并行环境下的全局通信，而不是单机小规模运行效率。本文结果也符合这一点：在当前单机原型规模内，PPPM仍然非常快。

随机批量方法的理论吸引力在于用随机估计降低每步交互规模\cite{jin2020random}。本文实现的RBSOG继承了这一思想，但实验表明，随机估计的计算收益只有在采样路径足够高效时才可能显现。否则，随机索引生成、样本过滤和无偏权重修正本身会吞噬理论节省。

高斯和分解提供了处理奇异核的一种光滑近似方式\cite{guo2025sum}。本文的direct reference实验说明，这条路径可以运行并能被系统评测，但当前核拟合策略还不够稳定。特别是 $M$ 增大导致误差变差这一现象，说明SOG分解在分子动力学静电场景中的参数化仍需专门研究。

\section{局限性}

本文仍有若干局限。

首先，体系是膜启发代理模型，不是全原子DPPC或真实细胞环境。因此，结论主要是算法工程层面的，不能直接外推为生物物理结论。

其次，当前实现仍是Python研究原型。尽管包含Numba加速，但与成熟C++/CUDA分子动力学引擎相比仍有明显差距。因此，本文比较的是原型实现而非生产级实现。

再次，规模扫描最大到 $N=1280$，尚不足以检验RBSOG在超大规模并行环境中的低通信潜力。当前实验只说明它在可运行单机规模内尚未体现优势。

最后，direct reference精度实验虽然覆盖了 $N=160,320,640$，但仍只使用有限快照和有限参数组合。未来需要更系统地区分SOG核误差、随机采样误差、短程截断误差和积分反馈误差。

\section{后续工作}

基于本文结果，后续工作应优先从以下方向展开：

（1）重构长程采样路径，避免生成大量随后被过滤的无效样本；

（2）重新设计SOG核拟合流程，比较不同拟合区间、权重约束和误差目标；

（3）增加更大规模和更长轨迹实验，区分短期压力噪声与长期结构趋势；

（4）开发多线程或GPU版本，真正测试RBSOG规避FFT通信的潜在优势；

（5）迁移到更真实的粗粒化或全原子体系，同时保持当前自动化实验与论文资产链路。

总体而言，当前工作完成了从方法想法到可复现实验体系的闭环，也明确指出了路线的真实困难。这比只报告单次有利结果更有价值，因为它为后续优化提供了可验证的靶点。
